%%%%%%%%%%%%%%%%%%%%%%%%%%%%%%%%%%%%%%%%%%%%%%%%%%%%%%%%%%%%%%%%%%%%%%
% Lớp 11 | Đạo hàm | Đạo hàm cấp hai
% Dạng 1D7?3-2: Đẳng thức có 𝑦và (𝑦′, 𝑦′′)
% Số câu: 20  (sinh tự động bởi tach_cau.py; từ nay thầy sửa trực tiếp file này)
%%%%%%%%%%%%%%%%%%%%%%%%%%%%%%%%%%%%%%%%%%%%%%%%%%%%%%%%%%%%%%%%%%%%%%

%%%%%%%%%%%%%%%% TRẮC NGHIỆM 4 PHƯƠNG ÁN %%%%%%%%%%%%%%%%

% nguồn: Ok_Toan11_TracNghiem_Lan2
\begin{ex}%[1D7H3-2]%[ma:cb769e01]
	 Cho hàm số $y=x\cos x$. Tìm hệ thức đúng trong các hệ thức sau.
	 \choice
		{$y''+y=\sin x+2x\cos x$} 
		 {$y''+y=2\sin x$}
		{$y''+y=-\sin  x+x\cos x$}
		{\True  $y''+y=-2\sin x$}
	\loigiai{ 
		Ta có $y'=\cos  x-x\sin x$ và $y''=-2\sin x-x\cos x$.\\
		Khi  đó $y''+y=-2\sin x-x\cos x+x\cos x=-2\sin x$.
	 }
\end{ex}

% nguồn: Ok_Toan11_TracNghiem_Lan2
\begin{ex}%[1D7H3-2]%[ma:f966f15b]
	 Cho hàm số $y=\sin^3x$. Khẳng định nào sau đây đúng?
	\choice 
		 {$y''+9y-\sin x=0$}
		{\True $y''+9y-6\sin x=0$} 
		{$y''+9y-6\cos  x=0$}
		{$y''+9y+6\sin x=0$} 
	 \loigiai{
		Ta có $y'=3\sin^2x\cdot\cos x$ và $y''=6\sin x\cdot\cos^2x-3\sin^3x$.\\ 
		Khi đó $y''+9y=6\sin  x\cdot\cos^2x-3\sin^3x+9\sin^3x=6\sin x\left(\sin^2x+\cos^2x\right)=6\sin x$.\\
		Vậy $y''+9y-6\sin x=0$. 
	} 
\end{ex}

% nguồn: Ok_Toan11_TracNghiem_Lan2
\begin{ex}%[1D7H3-2]%[ma:8da5df30]
	Cho hàm số $y=\sin 2x$. Mệnh đề nào sau đây đúng? 
	 \choice
		 {$y^2+\left(y'\right)^2=4$}
		{$y=y'\cdot\tan  2x$}
		{$4y-y''=0$} 
		 {\True $4y+y''=0$}
	\loigiai{ 
		Ta có  $y'=2\cos 2x$ và $y''=-4\sin 2x$.\\
		Vậy $4y+y''=4\sin 2x-4\sin 2x=0$. 
	 }
\end{ex}

% nguồn: Ok_Toan11_TracNghiem_Lan2
\begin{ex}%[1D7H3-2]%[ma:0ceb6412]
	 Cho hàm số $y=\sin 2x$. Hãy chọn câu đúng.
	\choice 
		 {$y^2+\left(y'\right)^2=4$}
		 {$4y-y''=0$}
		{\True $4y+y''=0$} 
		{$y=y'\tan 2x$} 
	 \loigiai{
		Ta có  $y'=2\cos 2x$ và $y''=-4\sin 2x$.\\
		 Vậy $4y+y''=4\sin 2x-4\sin 2x=0$.
	 }
\end{ex}

% nguồn: Ok_TracNghiemLop11_Lan1
\begin{ex}%[1D7H3-2]%[ma:52bc52b9]
	Cho hàm số $ y=x\cos x$. Tìm hệ thức đúng trong các hệ thức sau 
	\choice 
		{$y''+y=\sin x+2x\cos x$} 
		{$y''+y=2\sin x$} 
		{$y''+y=-\sin  x+x\cos x$}
		{\True  $y''+y=-2\sin x$}
	\loigiai{ 
		Ta có $y'=\text{cos}x-x\sin x\Rightarrow  y''=-2\sin x-x\cos x$.\\
	Khi đó $y''+y=-2\sin x-x\cos x+x\cos x=-2\sin x$.} 
\end{ex}

% nguồn: Ok_TracNghiemLop11_Lan1
\begin{ex}%[1D7H3-2]%[ma:cf5ad10b]
	Cho $ y=\sqrt{2x-x^2}$, tính giá  trị biểu thức $ A=y^3\cdot y''$.
	 \choice
		{$ 1$} 
		{$  0$}
		{\True $-1$} 
		{$4$} 
	 \loigiai{
		Ta có $  y'=\dfrac{1-x}{\sqrt{2x-x^2}}$, $y''=\dfrac{-1}{\left(\sqrt{2x-x^2}\right)^3}$.\\
	Do đó $  A=y^3.y''=-1$.}
\end{ex}

% nguồn: Ok_TracNghiemLop11_Lan1
\begin{ex}%[1D7H3-2]%[ma:84ac5e75]
	Cho hàm số $y = x \sin x$. Hệ thức nào sau đây  đúng?
	\choice 
		{$y'' - y =  2 \sin x$}
		{$y'' + y =  2 \sin x$}
		{$y'' -  y = 2 \cos x$}
		{\True $y''  + y = 2 \cos x$}
	\loigiai{ 
		 Ta có: $y' = \sin x + x \cos x$. Và
		$y''= \cos x + \cos x - x \sin x =  2 \cos x - x \sin x$. \\
		Vậy $y'' + y = 2 \cos x - x \sin x + x \sin x = 2 \cos x$. 
	 }
\end{ex}

% nguồn: Ok_Toan11_TracNghiem_Lan2
\begin{ex}%[1D7V3-2]%[ma:84f75293]
	Cho hàm số $y=\sqrt{1+3x-x^2}$.  Khẳng định nào dưới đây đúng?
	\choice 
		{\True  $\left(y'\right)^2+y\cdot y''=-1$}
		{$\left(y'\right)^2+2y\cdot y''=1$} 
		{$y\cdot y''-\left(y'\right)^2=1$} 
		 {$\left(y'\right)^2+y\cdot y''=1$}
	 \loigiai{
		Ta có $y=\sqrt{1+3x-x^2} \Rightarrow y^2=1+3x-x^2\Rightarrow 2y\cdot  y'=3-2x$.\\
		Suy ra $2\cdot\left(y'\right)^2+2y\cdot y''=-2  \Rightarrow\left(y'\right)^2+y\cdot y''=-1$.
	} 
\end{ex}

% nguồn: Ok_Toan11_TracNghiem_Lan2
\begin{ex}%[1D7V3-2]%[ma:8d1c9ecb]
	Cho $y=\sqrt{2x-x^2}$, tính giá  trị biểu thức $A=y^3\cdot y''$.
	\choice 
		 {$1$}
		{$0$} 
		 {\True $-1$}
		 {$2$}
	 \loigiai{
		Ta có $y'=\dfrac{1-x}{\sqrt{2x-x^2}}$  và $y''=\dfrac{-1}{\left(\sqrt{2x-x^2}\right)^3}$.\\
		 Vậy $A=y^3\cdot y''=\left(\sqrt{2x-x^2}\right)^3\cdot\dfrac{-1}{\left(\sqrt{2x-x^2}\right)^3}=-1$.
	} 
\end{ex}

% nguồn: Ok_TracNghiemLop11_Lan1
\begin{ex}%[1D7V3-2]%[ma:727ec089]
	Tất cả  các nghiệm của phương trình $y''=\sqrt{2}$ với $y=-\sqrt{3}\sin x-\cos x$ là
	 \choice
		{$x=\dfrac{\pi}{12}+k\pi \textrm{  hoặc }x=\dfrac{7\pi}{12}+k\pi \textrm{ với } k\in \mathbb{Z}$}
		{\True $x=\dfrac{\pi}{12}+k2\pi  \textrm{ hoặc }x=\dfrac{7\pi}{12}+k2\pi \textrm{ với } k\in \mathbb{Z}$}
		{$x=-\dfrac{\pi}{12}+k\pi  \textrm{ hoặc }x=\dfrac{5\pi}{12}+k\pi \textrm{ với } k\in \mathbb{Z}$}
		{$x=-\dfrac{\pi}{12}+k2\pi \textrm{ hoặc }x=\dfrac{5\pi}{12}+k2\pi \textrm{ với } k\in  \mathbb{Z}$}
	\loigiai{$y'=-\sqrt3\cos x+\sin  x, y''=\sqrt3 \sin x+\cos x=2\sin(x+\dfrac{\pi}{6}).$ \\
		Giải phương trình $2\sin(x+\dfrac{\pi}{6})=\sqrt2$ ta được nghiệm là $x=\dfrac{\pi}{12}+k2\pi \textrm{  hoặc }x=\dfrac{7\pi}{12}+k2\pi \textrm{ với } k\in \mathbb{Z}.$
	} 
\end{ex}

% nguồn: Ok_TracNghiemLop11_Lan1
\begin{ex}%[1D7B3-2]%[ma:e740401f]
	Cho hàm số $y = \dfrac{1}{{1 -  x}}$. Chọn khẳng định đúng trong các khẳng định sau:
	 \choice
		 {$y'' + y^3 = 0$}
		 {$y'' - y^3 = 0$}
		{$y'' +  2y^3 = 0$}
		{\True $y'' - 2y^3  = 0$}
	 \loigiai{
		Ta có  $y'=\dfrac{1}{(1-x)^2},\ y''=\dfrac{2}{(1-x)^3}$. Suy ra $y''-2y^3=0$.
	 }
\end{ex}


%%%%%%%%%%%%%%%% ĐÚNG/SAI %%%%%%%%%%%%%%%%

% nguồn: Ok_dungsai_lop11_Lan1
\begin{ex}%[1D7H3-2]%[ma:cecec529]
 \%[Dự án đề kiểm tra  Toán khối 11 GHKII NH23-24-Dot 1- Tên Huỳnh Kim Triển]\%[Deso 9- Sach CTST]
 Tính được đạo hàm cấp  hai của các hàm số sau. Khi đó:
  \choiceTF
  { $y=x^3-x^2+9x-5$ có $y''(-2)=14$} 
   {$y=2\cos \left(3x-\dfrac{\pi}{4}\right)$ có $y''\left(\dfrac{\pi}{6}\right)=9\sqrt{2}$}
  {\True $y=2\mathrm{e}^{2x-1}$ có  $y''(1)=8\mathrm{e}$}
  {\True $y=\ln (1-2x)$ có $y''(-3)=\dfrac{-4}{49}$} 
 \loigiai{\begin{itemize}
   \item Ta có
  \[y'=3 x^2-2 x+9 \Rightarrow y''=6 x-2.\] 
  Vậy  $y''(-2)=-14$.
   \item Ta có
   \[y'=-6 \sin \left(3 x-\dfrac{\pi}{4}\right) \Rightarrow y''=-18 \cos \left(3 x-\dfrac{\pi}{4}\right).\]
   Vậy $y''\left(\dfrac{\pi}{6}\right)=-9 \sqrt{2}$.
  \item Ta  có
  \[y'=4 \mathrm{e}^{2 x-1} \Rightarrow  y''=8 \mathrm{e}^{2 x-1}\]
  Vậy  $y^{\prime \prime}(1)=8 \mathrm{e}$.
  \item Ta có\[y'=\dfrac{(1-2 x)'}{1-2 x}=-\dfrac{2}{1-2  x} \Rightarrow y''=2\cdot \dfrac{(1-2 x)'}{(1-2 x)^2}=-\dfrac{4}{(1-2 x)^2}.\]
  Vậy $y''(-3)=\dfrac{-4}{49}$. 
\end{itemize}}
\end{ex}

% nguồn: Ok_dungsai_lop11_Lan1
\begin{ex}%[1D7H3-2]%[ma:6754f82b]
 Xét tính đúng, sai của các mệnh đề  sau
 \choiceTF 
  {\True Cho  hàm số $f(x)=x^3-3x^2+4x-6$. Nghiệm của bất phương trình $f''(x)\leqslant f'(x)-1$ là $x\in \left(-\infty;1\right]\cup \left[3;+\infty \right)$}
  {Cho hai hàm số $f(x)=x^4-4x^2+3$ và $g(x)=3+10x-7x^2$. Số nghiệm của phương trình $f''(x)+g'(x)=0$ là $1$} 
  {\True Cho  hàm số $y=3x^5-5x^4+3x-2$. Bất phương trình $y''<0$ có nghiệm là $x\in (-\infty;1)\setminus \{0\}$}
   {Cho hàm số $y=\sin 5x$. Tỉ số $\dfrac{y''}{y}=1$}
 \loigiai{ 
  \begin{itemchoice}
   \itemch $f(x)=x^3-3x^2+4x-6 \Rightarrow  f'(x)=3x^2-6x;f''(x)=6x-6$.\\
   $f''(x)\leqslant f'(x)-1\Leftrightarrow 6x-6\leqslant 3x^2-6x+4-1\Leftrightarrow 3x^2-12x+9\geqslant 0\Leftrightarrow \hoac{&x\leqslant  1\\&x\geqslant 3.} $
    \itemch $f'(x)=4x^3-8x$, $f''(x)=12x^2-8$, $g'(x)=10-14x$.\\
   $f''(x)+g'(x)=0 \Leftrightarrow 12x^2-8+10-14x=0 \Leftrightarrow 12x^2-14x+2=0 \Leftrightarrow  \hoac{&x=1\\&x=\dfrac{1}{6}.} $\\
   Vậy phương trình có $2$  nghiệm.
   \itemch $y'=15x^4-20x^3+3$, $y''=60x^3-60x^2$.\\ 
   $y''<0 \Leftrightarrow 60x^2(x-1)<0 \Leftrightarrow \heva{&x<1\\&x\ne 0.} $ 
   \itemch $y'=5\cos 5x$, $y''=-25\sin 5x\Rightarrow  \dfrac{y''}{y}=-25$.
  \end{itemchoice}
  }
\end{ex}

% nguồn: Ok_dungsai_lop11_Lan1
\begin{ex}%[1D7H3-2]%[ma:de20ae78]
 Cho hàm số $y=\sqrt{2x-x^2}$. Xét tính đúng sai của các mệnh đề sau 
 \choiceTF 
  {\True $y'=\dfrac{1-x}{\sqrt{2x-x^2}}$} 
  {\True  $y''=\dfrac{-1}{\left(\sqrt{2x-x^2}\right)^3}$}
   {$y^3\cdot y''=1$}
  {\True $y'+\left(x^3-2x^2\right)\cdot  y''=-\dfrac{1}{y}$}
 \loigiai{ 
  \begin{itemchoice}
   \itemch Đúng. Ta có $y'=\left(\sqrt{2x-x^2}\right)'\Leftrightarrow y'=\dfrac{1}{2\sqrt{2x-x^2}}\cdot\left(2x-x^2\right)'=\dfrac{1-x}{\sqrt{2x-x^2}}$. 
   \itemch Đúng vì 
   \begin{eqnarray*}
    &y''&=\dfrac{(1-x)'\cdot\sqrt{2x-x^2}-\left(\sqrt{2x-x^2}\right)'\cdot  (1-x)}{\left(\sqrt{2x-x^2}\right)^2}\\
    &&=\dfrac{-\sqrt{2x-x^2}-\dfrac{1-x}{\sqrt{2x-x^2}}\cdot  (1-x)}{\left(\sqrt{2x-x^2}\right)^2}\\
    &&=\dfrac{-\left(2x-x^2\right)-\left(1-x\right)^2}{\sqrt{2x-x^2}\cdot\left(\sqrt{2x-x^2}\right)^2}\\ 
     &&=\dfrac{-1}{\left(\sqrt{2x-x^2}\right)^3}.
   \end{eqnarray*}
   \itemch Sai vì $y^3\cdot  y''=\left(\sqrt{2x-x^2}\right)^3\cdot\dfrac{-1}{\left(\sqrt{2x-x^2}\right)^3}=-1$.
    \itemch Đúng vì $$y'+\left(x^3+2x^2\right)\cdot y''=\dfrac{1-x}{\sqrt{2x-x^2}}+\left(x^3-2x^2\right)\cdot\dfrac{-1}{\left(\sqrt{2x-x^2}\right)^3}=-\dfrac{1}{\sqrt{2x-x^2}}=-\dfrac{1}{y}.$$
  \end{itemchoice}
 } 
\end{ex}

% nguồn: Ok_dungsai_lop11_Lan1
\begin{ex}%[1D7H3-2]%[ma:4539527e]
 Xét tính đúng, sai của  các mệnh đề sau
 \choiceTF 
  {Cho hàm số $y=\sin 2x$. Khi  đó $4y+y''=2$}
  {\True Cho hàm số $y=x\cdot \sin x$. Ta luôn có $xy''-2y'+xy=-2\sin  x$}
  {\True Cho hàm số $y=x\cos2x$. Khi đó  $xy''+2(\cos 2x-y')+4xy=0$}
  {Cho hàm số $f\left(x\right)=\cos x+3x^{2}$. Số  nghiệm của phương trình $f(x)+f''(x)=0$ là $2$}
  \loigiai{
  \begin{itemchoice}
   \itemch  $y=\sin 2x \Rightarrow y'=2\cos 2x$.\\
   Ta có $y''=-4\sin  2x=-4y \Leftrightarrow 4y+y''=0$.
   \itemch Ta có  $y'=\sin x+x\cos x\Rightarrow y''=2\cos x-x\sin x=2\cos x-y$.\\
   Khi đó $xy''-2y'+xy=x(2\cos x-y)-2(\sin x+x\cos x)+xy=-2\sin x$. 
   \itemch $y=x\cos2x$;  $y'=(x\cos2x)’=\cos2x-2x\sin 2x$.\\
   $y''=(\cos2x-2x\sin 2x)’ =-2\sin 2x-2(\sin 2x+2x\cos2x) =-4\sin 2x-4x\cos2x$.\\ 
   Ta có  $xy''+2(\cos2x-y')+4xy =-4x\sin 2x-4x^2\cos2x+4x\sin 2x+4x^2\cos2x=0$.
   \itemch Ta có $f'\left(x\right)=-\sin x+6x$ và  $f''\left(x\right)=-\cos x+6$.\\
    $\Rightarrow f(x)+f''(x)=0 \Leftrightarrow 3x^2+6=0$ (vô nghiệm).
  \end{itemchoice}
  }
\end{ex}

% nguồn: Ok_dungsai_lop11_Lan1
\begin{ex}%[1D7H3-2]%[ma:86b32d73]
 Cho hàm số $y=16 \cos x +17 \sin x$.  Xét tính đúng sai của các mệnh đề sau
  \choiceTF
   {\True $y\left(\dfrac{\pi}{4}\right)=\dfrac{33\sqrt{2}}{2}$}
  {$y'=16\sin x+17\cos  x$}
  {\True $y''=-16\cos  x-17\sin x$}
  {$y''+y=1$} 
  \loigiai{
  \begin{itemchoice}
   \itemch  Ta có $y\left(\dfrac{\pi}{4}\right)=17\cos\dfrac{\pi}{4}+16\sin\dfrac{\pi}{4}=\dfrac{33\sqrt{2}}{2}$.
   \itemch $y'=-16 \sin x +17 \cos x$. 
   \itemch $y''=-16 \cos x -17  \sin x$.
   \itemch $y''+y=16 \cos x +17 \sin x -16  \cos x -17 \sin x =0$.
  \end{itemchoice}
 } 
\end{ex}

% nguồn: Ok_dungsai_lop11_Lan2
\begin{ex}%[1D7H3-2]%[ma:8d175eca]
	Cho hàm số $f(x)=\dfrac{1}{3x+2}$. 
	\choiceTF 
		{Tập xác định của hàm số là $\mathscr{D}=\left(\dfrac{2}{3};+\infty\right)$} 
		{\True Đạo hàm cấp một của hàm số là $f'(x)=\dfrac{-3}{(3x+2)^2}$} 
		{Đạo hàm cấp  hai của hàm số là $f''(x)=\dfrac{9}{(3x+2)^3}$}
		{\True  $\dfrac{f''(x)}{f(x)}+6f'(x)=0$}
	 \loigiai{
		\begin{itemchoice}
			\itemch Tập  xác định $\mathscr{D}=\mathbb{R}\setminus\left\{-\dfrac{2}{3}\right\}$.
			\itemch Ta  có $f(x)=(3x+2)^{-1}\Rightarrow f'(x)=-1\cdot (3x+2)'(3x+2)^{-2}=\dfrac{-3}{(3x+2)^2}$.
			\itemch Ta có $f'(x)=-3\cdot (3x+2)^{-2}\Rightarrow  f''(x)=(-3)\cdot (-2)\cdot (3x+2)'(3x+2)^{-3}=\dfrac{18}{(3x+2)^3}$.
			\itemch  Ta có $\dfrac{f''(x)}{f(x)}+6f'(x)=\dfrac{18}{(3x+2)^2}-\dfrac{18}{(3x+2)^2}=0$.
		\end{itemchoice}
	} 
\end{ex}

% nguồn: Ok_dungsai_lop11_Lan2
\begin{ex}%[1D7H3-2]%[ma:84097659]
	 Cho hàm số $y=x\ln(x+1)$.
	 \choiceTF
		{Hàm số đã cho có tập  xác định là $[0;+\infty)$}
		 {$y'=\ln(x+1)$}
		{\True $y''=\dfrac{1}{x+1}+\dfrac{1}{(x+1)^2}$} 
		{\True  $y'+(x+1)y''-2=\ln(x+1)$}
	\loigiai{ 
		\begin{itemchoice}
			\itemch Hàm số xác định khi $x+1>0\Leftrightarrow  x>-1$.\\
			Vậy tập xác định của hàm số đã cho là  $\mathscr{D}=(-1;+\infty)$.
			\itemch Ta có  $y'=x'\cdot\ln(x+1)+x\cdot (\ln(x+1))'=\ln(x+1)+\dfrac{x}{x+1}=\ln(x+1)+1-\dfrac{1}{x+1}$.
			 \itemch Ta có $y''=\left(\ln(x+1)+1-\dfrac{1}{x+1}\right)'=\dfrac{1}{x+1}+\dfrac{1}{(x+1)^2}$.
			\itemch Ta có $y'+(x+1)y''-2=\ln(x+1)+1-\dfrac{1}{x+1}+1+\dfrac{1}{x+1}-2=\ln(x+1)$. 
		\end{itemchoice}
	 }
\end{ex}

% nguồn: Ok_dungsai_lop11_Lan2
\begin{ex}%[1D7H3-2]%[ma:915c7a8c]
	Cho  hàm số $y=(x-1)\mathrm{e}^x$.
	 \choiceTF
		{\True Hàm số  đã cho xác định trên $\mathbb{R}$}
		 {$y'=(x-1)\mathrm{e}^x$}
		{\True  $y''=y'+\mathrm{e}^x$}
		{Phương trình $\left(y'\right)^2-y''\cdot\mathrm{e}^x+2\mathrm{e}^{2x}=0$ có hai nghiệm  phân biệt}
	\loigiai{ 
		\begin{itemchoice}
			\itemch Hàm số  đã cho xác định trên $\mathbb{R}$.
			\itemch Ta  có $y'=(x-1)'\cdot\mathrm{e}^x+(x-1)\cdot (\mathrm{e}^x)'=\mathrm{e}^x+(x-1)\mathrm{e}^x=x\mathrm{e}^x$.
			\itemch Ta có $y''=x'\cdot\mathrm{e}^x+x\cdot (\mathrm{e}^x)'=\mathrm{e}^x+x\mathrm{e}^x=\mathrm{e}^x+y'=y'+\mathrm{e}^x$. 
			\itemch Ta có \begin{eqnarray*}
			\left(y'\right)^2-y''\cdot\mathrm{e}^x+2\mathrm{e}^{2x}=0 &\Leftrightarrow&  x^2\cdot\mathrm{e}^{2x}-(x+1)\cdot\mathrm{e}^x\cdot\mathrm{e}^x+2\mathrm{e}^{2x}=0\\ &\Leftrightarrow & \mathrm{e}^{2x}\left(x^2-x-1+2\right)=0\\ &\Leftrightarrow &x^2-x+1=0.
		\end{eqnarray*}
		Phương trình $x^2-x+1=0$ vô nghiệm nên phương trình $\left(y'\right)^2-y''\cdot\mathrm{e}^x+2\mathrm{e}^{2x}=0$  vô nghiệm.
	\end{itemchoice}
 }
\end{ex}

% nguồn: Ok_dungsai_lop11_Lan1
\begin{ex}%[1D7B3-2]%[ma:070b70e5]
 Cho hàm số $y = \sin ^2 x$. 
 \choiceTF 
  {\True  $y'= 2\sin x \cos x$}
   {$y''=-2\cos 2x$}
  {$(y')^2 + (y'')^2 = 1$} 
  {\True $(y')^2 + (1 - 2y)^2  = 1$}
  \loigiai{Ta có
  \begin{itemchoice}
   \itemch Đúng, vì $y' = 2\sin x \cdot  (\sin x)' = 2\sin x \cos x = \sin 2x$.
   \itemch  Sai, vì $y''=2\cos 2x$.
   \itemch Sai, vì $(y')^2 + (y'')^2 = \sin^2 2x + 4\cos^2 2x =  1+3\cos^2 2x \ne 1$.
    \itemch Đúng, vì $(y')^2 + (1-2y)^2 = \sin^2 2x + (1-2\sin^2 x)^2 = \sin^2 2x + \cos^2 2x = 1$.
  \end{itemchoice}
  }
\end{ex}
